\documentclass[12pt,a4paper]{article}
\usepackage[utf8]{inputenc}
\usepackage[spanish]{babel}
\usepackage{amsmath,amssymb,amsfonts}
\usepackage{geometry}
\geometry{margin=2.5cm}

\title{\textbf{Compendio de Fórmulas de Óptica}}
\author{Física Óptica y Fotónica}
\date{\today}

\begin{document}

\maketitle

\section*{Introducción}
La óptica es la rama de la física que analiza el comportamiento y propagación de la luz. Se divide en óptica geométrica (rayos) y óptica física (ondas y fotones). A continuación se muestran sus fórmulas principales con sus unidades fundamentales en el SI.

\section{Óptica Geométrica}

\subsection{Reflexión y Refracción}

\subsubsection*{1. Leyes de Refracción y Reflexión [$n, \theta$]}
\paragraph*{Unidades básicas del SI:}
\begin{itemize}
    \item Índice de refracción [$n$]: Adimensional $[1]$
    \item Ángulos [$\theta$]: Radianes $[\text{rad}]$ (adimensional $[1]$)
\end{itemize}
\paragraph*{Explicación:} Describen el cambio de dirección que experimenta un rayo de luz al incidir sobre la interfaz entre dos medios transparentes.
\paragraph*{Desarrollo y Variaciones:}
\begin{align*}
\text{Índice de refracción:} \quad & n = \frac{c}{v} \\
\text{Ley de reflexión:} \quad & \theta_i = \theta_r \\
\text{Ley de Snell (refracción):} \quad & n_1 \sin\theta_1 = n_2 \sin\theta_2 \\
\text{Ángulo límite (reflexión total interna):} \quad & \theta_c = \arcsin\left(\frac{n_2}{n_1}\right) \quad (n_1 > n_2)
\end{align*}
donde $c$ es la velocidad de la luz en el vacío y $v$ la velocidad de la luz en el medio.
\paragraph*{Autor:} Willebrord Snellius / René Descartes.

\subsection{Lentes y Espejos}

\subsubsection*{2. Ecuación de los Lentes Delgados y Espejos [$f, s, s', P$]}
\paragraph*{Unidades básicas del SI:}
\begin{itemize}
    \item Distancias focal, de objeto e imagen [$f, s, s'$]: $[\text{m}]$
    \item Potencia óptica [$P$]: $[\text{m}^{-1}]$
\end{itemize}
\paragraph*{Explicación:} Relaciona las distancias del objeto y de la imagen respecto al centro óptico de la lente o espejo con su distancia focal.
\paragraph*{Desarrollo y Variaciones:}
\begin{align*}
\text{Ecuación de Gauss / Descartes:} \quad & \frac{1}{f} = \frac{1}{s} + \frac{1}{s'} \\
\text{Potencia óptica (dioptrías):} \quad & P = \frac{1}{f} \\
\text{Aumento lateral ($M$):} \quad & M = \frac{y'}{y} = -\frac{s'}{s} \\
\text{Ecuación del tallador de lentes (en aire):} \quad & \frac{1}{f} = (n - 1) \left( \frac{1}{R_1} - \frac{1}{R_2} \right)
\end{align*}
donde $R_1$ y $R_2$ son los radios de curvatura de las superficies de la lente.
\paragraph*{Autor:} René Descartes / Carl Friedrich Gauss.

\section{Óptica Física y Cuántica}

\subsection{Radiación de Cuerpo Negro}

\subsubsection*{3. Leyes de Radiación [$E, \lambda_{\max}, M_e$]}
\paragraph*{Unidades básicas del SI:}
\begin{itemize}
    \item Longitud de onda [$\lambda$]: $[\text{m}]$
    \item Exmitancia radiante [$M_e$]: $[\text{kg}\cdot\text{s}^{-3}]$
\end{itemize}
\paragraph*{Explicación:} Describen la distribución espectral y la potencia emitida por un cuerpo negro perfecto en función de la temperatura.
\paragraph*{Desarrollo y Variaciones:}
\begin{align*}
\text{Ley de Planck (Radiancia Espectral):} \quad & B_\lambda(\lambda, T) = \frac{2 \pi h c^2}{\lambda^5 \left(e^{\frac{h c}{\lambda k_B T}} - 1\right)} \\
\text{Ley de Desplazamiento de Wien:} \quad & \lambda_{\max} T = b \quad (b \approx 2.8977 \times 10^{-3} \text{ m}\cdot\text{K}) \\
\text{Ley de Stefan-Boltzmann:} \quad & M_e = \sigma T^4
\end{align*}
donde $\sigma$ es la constante de Stefan-Boltzmann y $b$ la constante de Wien.
\paragraph*{Autor:} Max Planck / Wilhelm Wien / Josef Stefan / Ludwig Boltzmann.

\end{document}
